\textbf{Время:} резерв

\textbf{Тезис:} архивный R-inv — эффективный параметр одномерного анализа с установленными ограничениями проверки.

\textbf{Визуал:} большой оригинальный график R-inv; показать центральности и исходные ошибки.

\textbf{Текст:} По горизонтали растёт k-T, по вертикали — R-inv в фемтометрах. В более центральных классах значения в целом выше. У самого центрального класса видна убывающая зависимость, но она не обязана одинаково проявляться в каждом классе центральности и каждом интервале kT.

Этот параметр получен одномерной аппроксимацией по q-inv. Он не является одним из трёх LCMS-радиусов и не обязан равняться их среднему. Усреднение по направлениям и движение пары меняют связь с источником.

Есть существенная граница достоверности. В формулах архивных слайдов встречаются ошибки записи, а код расчёта недоступен. Метод вычисления исходных ошибок также не установлен. Поэтому численные радиусы при принятом определении q независимо не подтверждены. Показываем сохранённый результат выполненной тогда работы, а не новую верификацию его точности.

\textbf{Пауза:} для количественного физического вывода нужно вернуться к событиям, проверить аппроксимацию и неопределённости.

\textbf{Переход:} Что потребуется, чтобы такой модельный результат действительно проверить экспериментом?

\textbf{Источники и условия:} Балинов (2025), слайд 14. В формулах архива есть ошибки записи; код того расчёта недоступен. \(R_{\rm inv}\) не равен отдельным LCMS-радиусам и радиусу ядра.
